\documentclass[12pt,a4paper]{article}
\usepackage{xeCJK}
\usepackage{geometry}
\usepackage{enumitem}
\usepackage{hyperref}
\usepackage{fancyhdr}

\geometry{left=2.5cm,right=2.5cm,top=3cm,bottom=3cm}
\setCJKmainfont{Noto Serif CJK TC}
\setCJKsansfont{Noto Sans CJK TC}

\pagestyle{fancy}
\fancyhf{}
\fancyhead[L]{內科部臨床能力委員會會議記錄}
\fancyhead[R]{2025年10月31日}
\fancyfoot[C]{\thepage}

\title{國立臺灣大學醫學院附設醫院新竹分院\\內科部臨床能力委員會(CCC)\\會議記錄}
\author{}
\date{2025年10月31日 12:30}

\begin{document}

\maketitle

\section*{會議資訊}
\begin{itemize}[leftmargin=*]
    \item \textbf{會議時間:} 2025年10月31日 12:30
    \item \textbf{會議形式:} 實體新竹台大分院五樓第一會議室，併行線上會議(Microsoft Teams)
    \item \textbf{主持人:} 謝慕揚醫師(教學部)
    \item \textbf{記錄:} 辜品馨(秘書室)
\end{itemize}

\section*{出席人員}
\begin{itemize}[leftmargin=*]
    \item 賴超倫醫師 (內科部主任)
    \item 謝慕揚醫師 (教學部,主講人)
    \item 趙政漢醫師 (內科部腎臟內科)
    \item 胡文皓醫師 (內科部腸胃內科)
    \item 耿立達主任 (內科部胸腔內科)
    \item 辜品馨 (內科部)
    \item 其他內科部次專科主任代表(線上)
\end{itemize}

\section{會議目的與背景}

謝慕揚醫師說明,臨床能力委員會(Clinical Competency Committee, CCC)的執行一直被認為較為複雜,但實際上全台灣私立醫學中心的許多科部都在進行。謝醫師於去年10月調至教學部,發現教學部能夠看到全國性及全世界性的趨勢。

本次會議主要目的為:
\begin{enumerate}[leftmargin=*]
    \item 介紹CCC的好處與EPA(Entrustable Professional Activities)概念
    \item 分享2025年8月巴塞隆納AMEE(歐洲醫學教育學會)的心得
    \item 討論內科部CCC的建立與運作機制
\end{enumerate}

\section{CBME概念說明}

\subsection{Competency-Based Medical Education(CBME)}

謝醫師解釋目前的教育理念為「以臨床能力為基礎的評量機制」(CBME):

\begin{itemize}[leftmargin=*]
    \item \textbf{以時間為本(Time-based):} 現行制度多採此方式,住院醫師過了一年就自動晉升,除非被發現有問題。但即使有問題,仍有很大機會讓他晉升。
    \item \textbf{以能力為本(Competency-based):} 國內外全部推動此制度,因為有些人學習速度較快，本來就可以更快進行能力升等，學習下一階段的臨床技能。同時對學生學習也是肯定，重點是更能提前發現學生學習的障礙。
\end{itemize}

謝醫師舉例說明,過去確實有制度可以讓不適任者無法晉升,但這是少數案例。他提到與胡副主任同年時,有一位同事在急診出事,因此在升R4時被內科部主任決定不能晉升。

\subsection{配合ACGME六大核心能力}

委員會的核心理念配合美國ACGME(Accreditation Council for Graduate Medical Education)的六大核心能力架構。

\section{臨床能力委員會(CCC)架構}

\subsection{目前執行現況}

\begin{itemize}[leftmargin=*]
    \item \textbf{兒科部:} 已執行CCC,每月召開一次會議,為國內醫學中心中表現優異者
    \item \textbf{影像醫學部:} 住院醫師與台大總院合作評估臨床能力,在教學醫院評鑑中獲得高度評價
    \item \textbf{內科部:} 目前規劃建立CCC制度
\end{itemize}

\subsection{委員會組成}

\begin{itemize}[leftmargin=*]
    \item \textbf{主席:} 內科部主任(賴主任)
    \item \textbf{委員:} 各次專科主任代表
    \item \textbf{醫事人員代表:} 護理部同仁
    \item \textbf{外部委員:} 急診醫學部醫師(已由主任圈選)
    \item \textbf{行政支援:} 秘書室
\end{itemize}

\subsection{委員資格}

\begin{itemize}[leftmargin=*]
    \item 具備能力、品格
    \item 有帶學生的經驗
    \item 符合一定教學年資
    \item 遵守利益迴避原則(可提供意見但不參與涉及自身學生的表決)
\end{itemize}

\section{委員會職責與運作}

\subsection{主要職責}

\begin{enumerate}[leftmargin=*]
    \item 評估學生/住院醫師的臨床能力,決定是否具備晉升資格
    \item 針對未達標準者,決定是否需要補訓或其他處置
    \item 提供系統性的評估機制,取代過去個別主治醫師的單獨評價(「舉紅旗模式」)
    \item 分擔個別教師的評估壓力與責任
\end{enumerate}

\subsection{適用對象}

\begin{itemize}[leftmargin=*]
    \item PGY(Post-Graduate Year)住院醫師
    \item 醫學系學生
    \item 專科護理師(Nurse Practitioner)
    \item 未來可擴展至Fellow等其他培訓人員
\end{itemize}

\subsection{會議頻率建議}

\begin{itemize}[leftmargin=*]
    \item \textbf{PGY:} 建議每3個月評估一次
    \item \textbf{醫學生:} 2-3週評估一次
    \item 可整併至現有會議中進行,減少額外負擔
    \item 可召開臨時會議或定期會議
\end{itemize}

\section{EPA(Entrustable Professional Activities)制度}

\subsection{EPA概念說明}

EPA為「可委任的專業活動」,是根據機構與單位需求而制定的評估項目。EPA的優點:

\begin{itemize}[leftmargin=*]
    \item 可適用於不同職類:呼吸治療師、住院醫師、NP等
    \item 將臨床能力具體化為可評估的項目
    \item 設定明確的等級標準
\end{itemize}

\subsection{EPA等級制度}

以ACGME為基準,設定Level 1至Level 4的能力等級:
\begin{itemize}[leftmargin=*]
    \item \textbf{Level 1:} 觀察學習階段
    \item \textbf{Level 2:} 在督導下執行
    \item \textbf{Level 3:} 可獨立執行但需監督
    \item \textbf{Level 4:} 完全獨立執行
\end{itemize}

\subsection{各次專科EPA範例}

謝醫師展示已完成的EPA文件:

\subsubsection{胸腔科EPA(與張醫師、溫醫師、陳思宇醫師共同制定)}

\begin{enumerate}[leftmargin=*]
    \item 胸腔科病史理學檢查與風險評估
    \item 胸腔穿刺術與胸管置放
    \item 其他臨床技能
    \item 支氣管鏡檢查(EPS)
    \item 其他進階技術
\end{enumerate}

針對不同訓練階段設定不同要求:
\begin{itemize}[leftmargin=*]
    \item \textbf{醫學生:} EPS為NA(不適用),但胸腔穿刺可達Level 3
    \item \textbf{住院醫師:} 需達到更高等級
    \item \textbf{NP:} 有特定的EPA項目
\end{itemize}

\subsubsection{腸胃內科EPA}

共11項EPA,可依不同職類調整:
\begin{itemize}[leftmargin=*]
    \item PGY可能需完成1-5項
    \item NP需完成1-7項  
    \item 主治醫師應完成1-11項
\end{itemize}

\subsection{其他科部EPA現況}

謝醫師表示已花費3個月時間,完成內科部10個次專科的EPA文件,包括:
\begin{itemize}[leftmargin=*]
    \item 心臟血管科
    \item 胸腔科
    \item 腸胃內科
    \item 其他次專科(共約100個EPA項目)
\end{itemize}

\section{評估流程與方法}

\subsection{評估資料來源}

\begin{enumerate}[leftmargin=*]
    \item \textbf{E-portfolio:} 全醫院使用的網頁版學生學習歷程評估系統
    \item \textbf{服務量統計:}
    \begin{itemize}
        \item 侵入性技術執行次數(如CVP、A-line等,可從portal系統查詢)
        \item 收治病人數(可從院編查詢)
        \item 其他客觀數據
    \end{itemize}
    \item \textbf{Mini-CEX評分表:} 現行使用的1-9分評分系統
    \item \textbf{360度評估}
    \item \textbf{Case-based討論}
    \item \textbf{CPT執行記錄}
\end{enumerate}

\subsection{評估程序}

\begin{enumerate}[leftmargin=*]
    \item 各專科主任或熱心教學的主治醫師進行定期評估
    \item 將評估結果提交CCC會議討論
    \item 委員會審查學生/住院醫師的各項表現
    \item 針對特定EPA項目進行能力評定
    \item 投票表決是否通過晉升
\end{enumerate}

\subsection{表決方式(參考ACLS委員會模式)}

謝醫師分享ACLS師資培訓的CCC運作經驗:

\begin{itemize}[leftmargin=*]
    \item 12位委員(來自6個醫學會,每學會3位代表)
    \item 考生上台演示後,委員給予評分
    \item 能力評核時，由秘書唱名,委員舉手表決
    \item 需達七票(12票中)才通過
    \item 60位考生的CCC可在25分鐘內完成
    \item 最後確認是否有委員認為應「一定要抓下」的人選
    \item 詢問是否有需要「加分」的特殊情況
\end{itemize}

\section{現行問題與CCC的解決方案}

\subsection{目前面臨的挑戰}

賴主任補充說明目前的困境:

\begin{enumerate}[leftmargin=*]
    \item \textbf{評估壓力集中:} 過去幾年評分壓力都落在少數一兩位老師身上
    \item \textbf{教師承擔風險:} 
    \begin{itemize}
        \item 老師需單獨承擔給予不及格的壓力
        \item 曾有教師寫email評論學生表現,反而成為衛福部申訴時的唯一依據
        \item 醫院需要該email來向衛福部說明學生確實表現不佳
    \end{itemize}
    \item \textbf{教師拒絕帶學生:} 有些教師已開始拒絕帶有問題的學生
    \begin{itemize}
        \item 擔心投入大量臨床時間
        \item 擔心最後被家屬或學生投訴
        \item 曾有學生連續2個月有問題後,第3個月沒有老師願意帶
    \end{itemize}
    \item \textbf{PGY學生心態問題:}
    \begin{itemize}
        \item 部分PGY來醫院只是為了申請特定小科
        \item 若目標不一致,可能第一年就離職
        \item 有案例:10月公告PGY-2志願未錄取,11月就停止工作(僅做3個月就離開)
    \end{itemize}
    \item \textbf{擺爛學生處理困難:}
    \begin{itemize}
        \item 學生知道「滿一年就拿得到台大PGY資格」
        \item 持續擺爛,不聽指導
        \item 醫院收了PGY就要「奉陪到底」,無法讓其離開
    \end{itemize}
\end{enumerate}

\subsection{CCC制度的優勢}

\begin{enumerate}[leftmargin=*]
    \item \textbf{團體評估機制:} 不再由單一主管或教師承擔決定權
    \item \textbf{公平公正:} 類似醫院的懲審考績委員會,通過委員會表決
    \item \textbf{分散壓力:} 多位委員共同分擔評估責任
    \item \textbf{客觀評估:} 
    \begin{itemize}
        \item 有presenting teacher陳述觀察
        \item 有客觀的360度評估
        \item 有多面向的評分
    \end{itemize}
    \item \textbf{明確標準:} 透過EPA設定清楚的能力標準
    \item \textbf{申訴權利:} 學生仍保有申訴管道
\end{itemize}

\section{未來規劃與建議}

\subsection{謝醫師提案}

\begin{enumerate}[leftmargin=*]
    \item \textbf{畢業標準:} 須達到所有核心EPA的Level 4(獨立執行)
    % \item \textbf{未達標準者的處理:}
    % \begin{itemize}
    %     \item 若持續未達畢業門檻,讓其留下來
    %     \item 僅負責基本工作:接電話、開order、打病歷
    %     \item 在臨床能力委員會認定下,只能執行基本業務
    %     \item 除非完成更多訓練,否則無法畢業
    %\end{itemize}
    \item \textbf{補訓或退訓機制:}
    \begin{itemize}
        \item 賴主任建議:採用「及格/可補訓/退訓」三級制
        \item 補訓:給予改善機會(愛的教育)
        \item 補訓太多次則退訓
        \item 退訓後可轉至其他醫院
        \item 不提供不合格的certificate
    \end{itemize}
\end{enumerate}

\subsection{會議整併建議}

\begin{itemize}[leftmargin=*]
    \item 可配合現有的PGY座談會時間
    \item 將會議時間切割為兩個議題
    \item 減少額外會議負擔
    \item 確保委員都能出席
\end{itemize}

\subsection{不同職類的EPA規劃}

委員討論認為應針對不同職類分別制定EPA文件:

\begin{itemize}[leftmargin=*]
    \item \textbf{住院醫師:} 獨立的EPA文件
    \item \textbf{PGY:} 內科部PGY教學文件
    \item \textbf{醫學生:} 內科部醫學生教學文件
    \item \textbf{專科護理師:} 內科部NP教學文件
    \item \textbf{Fellow:} 另外制定
\end{itemize}

不建議將所有職類放在同一份文件的不同部分,以免內容過於複雜。

\subsection{各科EPA數量建議}

依據學生/住院醫師所在單位,設定需完成的EPA數量:
\begin{itemize}[leftmargin=*]
    \item 若在腸胃病房,需完成數個GI的EPA
    \item 若在CV病房,需完成數個CV的EPA
    \item 不要求跨科完成(如在腸胃科不需做chest或nephro的EPA)
    \item 不同層級有不同要求:
    \begin{itemize}
        \item 住院醫師:可能需完成5-7個
        \item PGY:3-5個
        \item 醫學生:3個
        \item NP:3-5個(但內容不同)
    \end{itemize}
\end{itemize}

\section{技術支援與工具}

\subsection{AI輔助工具應用}

謝醫師說明已運用AI技術協助EPA制定與會議記錄:

\begin{enumerate}[leftmargin=*]
    \item \textbf{EPA文件生成:}
    \begin{itemize}
        \item 使用AI協助,3個月完成10個次專科的EPA
        \item 單靠人力無法在短時間完成
    \end{itemize}
    
    \item \textbf{評估系統:}
    \begin{itemize}
        \item 使用Python撰寫HTML程式
        \item 可處理Excel檔案中的評分資料
        \item 自動生成所需報表
        \item 平均10分鐘可完成一位學生的EPA評估
    \end{itemize}
    
    \item \textbf{會議記錄:}
    \begin{itemize}
        \item 使用Teams錄音並自動生成逐字稿
        \item 將逐字稿輸入AI處理
        \item 可在15分鐘內產出會議記錄
        \item 謝醫師承諾當晚就將本次會議記錄整理完成給主任
    \end{itemize}
    
    \item \textbf{AI訓練方式:}
    \begin{itemize}
        \item 將過去的教學紀錄、讀書會、口頭報告等文件輸入
        \item 包含個人的想法思路與見解
        \item 最成功的案例:將過去出的題目全部輸入
        \item AI可記住使用者的邏輯與要求
        \item 使用ChatGPT、Claude或Gemini等工具(每月約500元)
    \end{itemize}
\end{enumerate}

\subsection{現有系統與資源}

\begin{itemize}[leftmargin=*]
    \item \textbf{E-portfolio:} 全院統一使用的學習歷程系統
    \item \textbf{Portal系統:} 可查詢侵入性技術執行記錄、收治病人數等
    \item \textbf{Mini-CEX:} 現行評分表(1-9分制)
    \item \textbf{教學部支援:} 教學部的秀成將協助建置相關系統
\end{itemize}

\section{下次會議規劃}

\subsection{會議目標}

\begin{enumerate}[leftmargin=*]
    \item 示範實際CCC運作流程
    \item 展示AI輔助的評估系統
    \item 進行第一次實際的CCC評估
    \item 完成病情解釋等EPA的錄音與評估流程示範
\end{enumerate}

\subsection{預期效益}

\begin{itemize}[leftmargin=*]
    \item 內科部教學評鑑可順利通過
    \item 跟上其他科部(影像醫學部、兒科部、婦產部)的進度
    \item 護理部也已開始執行類似制度
    \item 避免在教學評鑑時被委員指正
\end{itemize}

\section{會議決議}

\begin{enumerate}[leftmargin=*]
    \item \textbf{同意建立內科部臨床能力委員會(CCC)}
    \item \textbf{委員會組成:}
    \begin{itemize}
        \item 主席:內科部主任
        \item 委員:各次專科主任代表
        \item 醫事人員代表
        \item 外部委員:急診醫學部醫師
    \end{itemize}
    \item \textbf{採用EPA制度進行能力評估}
    \item \textbf{會議整併至現有PGY座談會時段}
    \item \textbf{針對不同職類分別制定EPA文件}
    \item \textbf{下次會議將進行實際CCC評估示範}
\end{enumerate}

\section{散會}

會議於討論AI輔助系統與下次會議規劃後結束。

\vspace{1cm}

\noindent
\textbf{記錄人員:} 辜品馨

\noindent
\textbf{會議記錄完成日期:} 2025年10月31日

\end{document}